\documentclass[11pt]{letter}
\usepackage[utf8]{inputenc}
\usepackage[T1]{fontenc}
\usepackage[margin=2cm]{geometry}
\usepackage{url}

\setlength{\parskip}{4pt}
\longindentation=0pt

\signature{Magnus Lodefalk}
\address{\"Orebro University School of Business\\SE-701 82 \"Orebro, Sweden\\magnus.lodefalk@oru.se}

\begin{document}

\begin{letter}{The Editor\\Economics Letters}

\opening{Dear Editor,}

Enclosed for your consideration is an original research paper entitled ``Same Storm, Different Boats: Generative AI and the Age Gradient in Hiring''.

The paper asks whether the post-2022 divergence between stock prices and job postings reflects generative AI or monetary tightening, and what each channel implies for who is being hired. We exploit the seven-month gap between the Riksbank's first rate hike (April 2022) and the ChatGPT launch (November 2022) as a within-Sweden timing test, combining the full population of Swedish job advertisements (2020--2026) with monthly employer-declaration register data covering all employed individuals (2019--2025) in an employer-level difference-in-differences design with employer$\times$month fixed effects.

The aggregate posting decline tracks monetary tightening. Within employers, however, employment of workers aged 22--25 in AI-exposed occupations falls 5.5 per cent by mid-2025, concentrated among young women, mainly through reduced hiring, and accelerating only after ChatGPT. Population-level employment remains stable, so the pattern is one of reallocation rather than net destruction. The contribution is this reconciliation: aggregate vacancy declines reflect monetary tightening, while AI reshapes hiring composition inside firms.

The body is 1{,}956 words (within the 2{,}000 limit) with a 91-word abstract and an Online Appendix. The work is original and unpublished, and all authors approve the submission.

\closing{Yours sincerely,}

\end{letter}
\end{document}
